% Elsevier 3p Format - Two Column
% Based on Spatial-CvM Proofs Draft
\documentclass[final,3p,twocolumn,times,authoryear]{elsarticle}

% Packages allowed by Elsevier
\usepackage{amsmath,amssymb,amsthm}
\usepackage{mathtools}
\usepackage{bm}
\usepackage{booktabs}
\usepackage{hyperref}
\usepackage{cleveref}

% Theorem environments
\theoremstyle{definition}
\newtheorem{definition}{Definition}[section]
\newtheorem{theorem}{Theorem}[section]
\newtheorem{lemma}{Lemma}[section]
\newtheorem{proposition}{Proposition}[section]
\newtheorem{remark}{Remark}[section]

% Journal information
\journal{Journal of Statistical Planning and Inference}

\begin{document}

\begin{frontmatter}

\title{Mathematical Proofs for the Spatial Cramér--von Mises Test:\\Formalization in Lean 4}

\author[1]{Spatial-CvM Formalization Project}
\ead{project@spatialcvm.org}

\affiliation[1]{organization={Department of Statistics},
             addressline={University Research Center},
             city={City},
             postcode={12345},
             country={Country}}

\begin{abstract}
This paper presents the mathematical foundations of the Spatial Cramér--von Mises (CvM) test for assessing homogeneity of spatially dependent random fields. We establish asymptotic theory under $\alpha$-mixing dependence, including weak convergence of kernel-smoothed empirical processes to Gaussian limits, Mercer decomposition for chi-square approximations, and the functional delta method. All proofs are designed for formalization in Lean 4, with explicit tracking of axioms versus constructive results. The key innovation---fixed bandwidth kernel smoothing---ensures non-degenerate asymptotic distributions unlike classical shrinking-bandwidth theory.
\end{abstract}

\begin{keyword}
Spatial statistics \sep Cramér--von Mises test \sep empirical processes \sep weak convergence \sep formal verification \sep Lean 4
\MSC[2020]{Primary 62G10; Secondary 62M40, 60F17}
\end{keyword}

\end{frontmatter}

% ============================================
\section{Introduction}
\label{sec:intro}
% ============================================

The Spatial Cramér--von Mises test assesses homogeneity of spatially dependent random fields. Our formalization establishes three main theorems:

\begin{enumerate}
    \item[\textbf{T1}] Weak convergence of the kernel-smoothed empirical process to a Gaussian process in $\ell^\infty[0,1]$ (Section~\ref{sec:theorem1}).
    \item[\textbf{T2}] Joint convergence at multiple spatial locations and the limiting chi-square distribution (Section~\ref{sec:theorem2}).
    \item[\textbf{T3}] Functional delta method for the test statistic's asymptotic distribution (Section~\ref{sec:theorem3}).
\end{enumerate}

The mathematical framework builds on four pillars: Davydov's inequality for covariance bounds, Lindeberg CLT for finite-dimensional convergence, Arzelà--Ascoli for tightness, and Mercer's theorem for spectral decomposition. Our Lean 4 formalization tracks which components are proved versus axiomatized, providing a roadmap for completion.

% ============================================
\section{Preliminaries}
\label{sec:prelim}
% ============================================

\subsection{Location Space and Kernels}

\begin{definition}[Location Space]
Let $\mathcal{L} = \mathbb{R}^2$ be the location space with Euclidean metric $d(s_1, s_2) = \|s_1 - s_2\|_2$.
\end{definition}

\begin{definition}[Kernel Function]
\label{def:kernel}
A function $K: \mathbb{R} \to \mathbb{R}$ is a \emph{kernel} if:
\begin{enumerate}
    \item[\textbf{(K1)}] \textbf{Symmetry:} $K(x) = K(-x)$ for all $x \in \mathbb{R}$
    \item[\textbf{(K2)}] \textbf{Boundedness:} $|K(x)| \leq B$ for some $B > 0$
    \item[\textbf{(K3)}] \textbf{Compact Support:} $K(x) = 0$ for $|x| > r$
    \item[\textbf{(K4)}] \textbf{Lipschitz:} $|K(x) - K(y)| \leq L|x - y|$
    \item[\textbf{(K5)}] \textbf{Measurability:} $K$ is Borel measurable
\end{enumerate}
\end{definition}

\begin{definition}[Scaled Kernel]
For bandwidth $h > 0$:
\begin{equation}
K_h(x) = \frac{1}{h^2} K\left(\frac{x}{h}\right)
\end{equation}
\end{definition}

\subsection{Mixing and Dependence}

The spatial process satisfies $\alpha$-mixing with coefficients $\alpha(d)$ satisfying $\sum_{d=0}^\infty \alpha(d) < \infty$. Davydov's inequality (Davydov 1968; see also Rio 2013, Theorem~1.1) bounds covariances for centered indicator functions with a crucial exponent:
\begin{equation}
|\!\operatorname{Cov}(\mathbb{1}\{X_0 \leq x\} - F(x), \mathbb{1}\{X_d \leq y\} - F(y))\!| \leq 4 \cdot \alpha(d)^{\frac{\delta}{2+\delta}}
\end{equation}
where $\delta > 0$ is arbitrary. The Rio/Davydov exponent $\frac{\delta}{2+\delta} > 0$ ensures summability when $\alpha(d)$ decays geometrically, since $\sum_d \rho^{\frac{\delta d}{2+\delta}} < \infty$ for $\rho \in (0,1)$.

% ============================================
\section{Lemma 1: Asymptotic Covariance Structure}
\label{sec:lemma1}
% ============================================

\subsection{The Gamma Kernel}

\begin{definition}[Gamma Kernel]
\label{def:gamma-kernel}
For kernel $K$, bandwidth $h > 0$, lag $u \in \mathbb{R}$:
\begin{equation}
\gamma_{K,h}(u) = \int_{\mathbb{R}} K_h(v) K_h(v - u) \, dv
\end{equation}
\end{definition}

\begin{theorem}[Non-vanishing Variance]
\label{thm:nonvanish}
For any non-degenerate kernel $K$ and $h > 0$:
\begin{equation}
\gamma_{K,h}(0) = \int_{\mathbb{R}} K_h(v)^2 \, dv > 0
\end{equation}
\end{theorem}

\begin{proof}
Since $K$ is not identically zero, there exists a set of positive measure where $K \neq 0$. By (K2)--(K3), $K_h \in L^2(\mathbb{R})$. The integral of a non-negative, not-identically-zero function over positive measure is strictly positive.
\end{proof}

\begin{remark}[Key Innovation]
Theorem~\ref{thm:nonvanish} fails for shrinking bandwidth ($h \to 0$): then $\gamma_{K,h}(0) \to 0$. Fixed $h > 0$ ensures non-degenerate asymptotics.
\end{remark}

\subsection{The Gamma Operator}

\begin{definition}[Gamma Operator]
\label{def:gamma-op}
The covariance operator:
\begin{equation}
\Gamma_{K,h}(s_1, s_2) = \int_{\mathbb{R}} K_h(u - s_1) K_h(u - s_2) \, du
\end{equation}
\end{definition}

\begin{lemma}[Symmetry]
$\Gamma_{K,h}(s_1, s_2) = \Gamma_{K,h}(s_2, s_1)$
\end{lemma}

\begin{lemma}[Summability via Davydov-Rio]
\label{lem:summable-davydov}
Under geometric mixing $\alpha(d) \leq C\rho^d$ and the Rio/Davydov bound $|\gamma_d(y,z)| \leq 4 \cdot \alpha(d)^{\frac{\delta}{2+\delta}}$:
\begin{equation}
\label{eq:summability-bound}
\sum_{d=0}^\infty |\gamma_d(y,z)| \leq 4 \sum_{d=0}^\infty (C\rho^d)^{\frac{\delta}{2+\delta}} = \frac{4 C^{\frac{\delta}{2+\delta}}}{1 - \rho^{\frac{\delta}{2+\delta}}} < \infty
\end{equation}
\end{lemma}

\begin{proof}[Proof Sketch]
By Lemma~\ref{lem:lag-regroup} (Lag Regroup Identity), the variance decomposes by lag. The coefficient at lag $m$ is bounded by $2n$ when $|a_i| \leq 1$. Combining with the Rio/Davydov bound:
\begin{equation}
\Bigl|\sum_{i,j} a_i a_j \gamma(j-i)\Bigr| \leq 2n \sum_{|m| < n} |\gamma(m)| \leq 2n \cdot \frac{8 C^{\frac{\delta}{2+\delta}}}{1 - \rho^{\frac{\delta}{2+\delta}}}
\end{equation}
which is $O(n)$. The normalized variance is therefore $O(1/n) \cdot O(n) = O(1)$, as required.
\end{proof}

\begin{remark}[Key Exponent]
The exponent $\frac{\delta}{2+\delta}$ is \emph{sharp} for $L^p$ spaces: in Rio's framework with $p = 2+\delta$ and $q$ conjugate, the Davydov bound optimizes to this form. For summable $\alpha(d)$ and any $\delta > 0$, this gives sufficient decay.
\end{remark}

\begin{lemma}[Lag Regroup Identity]
\label{lem:lag-regroup}
For double sums over covariance structures:
\begin{equation}
\sum_{i,j=0}^{n-1} a_i a_j \gamma(j-i) = \sum_{m=-(n-1)}^{n-1} \gamma(m) \sum_{i: \text{valid}} a_i a_{i+m}
\end{equation}
where the inner sum has at most $n-|m|$ terms, giving coefficient bound $n-|m| \leq n$.
\end{lemma}

% ============================================
\section{Theorem 1: Weak Convergence}
\label{sec:theorem1}
% ============================================

\begin{theorem}[Main]
\label{thm:weak-conv}
Under Assumptions 1--4, the centered empirical process converges weakly in $\ell^\infty[0,1]$:
\begin{equation}
\sqrt{n}\bigl(\hat{H}_{n,h} - H_0\bigr) \stackrel{d}{\longrightarrow} \mathcal{G}
\end{equation}
where $\mathcal{G}$ is a zero-mean Gaussian process with covariance $\Gamma$.
\end{theorem}

\begin{proof}[Proof Sketch]
The proof follows standard empirical process theory via Prokhorov's theorem:

\textbf{Step 1: Finite-dimensional convergence.} Fix $0 \leq y_1 < \dots < y_k \leq 1$. The vector $\bigl(\sqrt{n}(\hat{H}_{n,h}(y_j) - H_0(y_j))\bigr)_{j=1}^k$ converges by Lindeberg CLT for $\alpha$-mixing arrays (El Machkouri--Volný--Wu, 2013).

\textbf{Step 2: Tightness.} Boundedness + equicontinuity via kernel Lipschitz property implies relative compactness by Arzelà--Ascoli.

\textbf{Step 3: Portmanteau.} FDD convergence + tightness $\Rightarrow$ weak convergence in $\ell^\infty[0,1]$.
\end{proof}

\begin{remark}[Formalization Status]
Theorem~\ref{thm:weak-conv} requires Prokhorov's theorem in non-separable spaces---not yet in Mathlib. Estimated completion: 1--2 years.
\end{remark}

% ============================================
\section{Theorem 2: Weighted Chi-Square Limit}
\label{sec:theorem2}
% ============================================

\begin{theorem}[Main]
\label{thm:chisq}
Under Assumptions 1--4, the test statistic converges:
\begin{equation}
T_n \stackrel{d}{\longrightarrow} \sum_{m=1}^\infty \lambda_m^* \cdot \chi^2_{K-1,m}
\end{equation}
where $\lambda_m^*$ are eigenvalues of the contrast covariance operator.
\end{theorem}

\begin{proof}[Proof Sketch]

\textbf{Step 1: Continuous Mapping.} $T_n = \Phi(\sqrt{n}(\hat{H}_{n,h} - H_0))$ where $\Phi(f) = \int_0^1 f(y)^2 \, dH_0(y)$ is continuous on $(\ell^\infty, \|\cdot\|_\infty)$.

\textbf{Step 2: Mercer Decomposition.} By Mercer's theorem:
\begin{equation}
\mathcal{G}(y) = \sum_{m=1}^\infty \sqrt{\lambda_m^*} \, \varphi_m^*(y) \, Z_m
\end{equation}
where $\varphi_m^*$ are eigenfunctions and $Z_m \sim N(0,1)$ i.i.d.

\textbf{Step 3: Spectral Representation.} Substituting:
\begin{equation}
\int_0^1 \mathcal{G}(y)^2 \, dH_0(y) = \sum_m \lambda_m^* Z_m^2 = \sum_m \lambda_m^* \chi^2_{1,m}
\end{equation}
The contrast subspace has dimension $K-1$, giving the final form.
\end{proof}

% ============================================
\section{Theorem 3: Functional Delta Method}
\label{sec:theorem3}
% ============================================

\begin{theorem}[Multivariate Extension]
\label{thm:delta}
For multivariate marks $\mathbf{Y}_i = (Y_{i,1}, \dots, Y_{i,p})$ with copula $C$:
\begin{equation}
T_n^{(p)} \stackrel{d}{\longrightarrow} \sum_{m=1}^\infty \lambda_m^{*,(p)} \cdot \chi^2_{K-1,m}
\end{equation}
\end{theorem}

\begin{proof}[Proof Sketch]
\textbf{Step 1: Sklar's Theorem.} Decompose via copula: $\mathbf{Y}_i = (F_1^{-1}(U_{i,1}), \dots, F_p^{-1}(U_{i,p}))$.

\textbf{Step 2: Hadamard Differentiability.} The quantile map $\Phi_p$ is Hadamard differentiable with derivative $1/f_j(F_j^{-1}(\cdot))$.

\textbf{Step 3: Chain Rule.} Apply Theorem~\ref{thm:weak-conv} to the copula process.
\end{proof}

% ============================================
\section{New Results: Algebraic Foundations}
\label{sec:new}
% ============================================

\subsection{Abel Summation (PROVED in Lean)}

\begin{theorem}[Abel Summation Identity]
For sequences $a_i, b_i$ and $n \geq 1$:
\begin{equation}
\sum_{i=0}^{n-1} a_i(b_{i+1} - b_i) = a_n b_n - a_0 b_0 - \sum_{i=0}^{n-1} (a_{i+1} - a_i) b_{i+1}
\end{equation}
\end{theorem}

\begin{proof}[Lean Implementation]
Induction on $n$ using telescoping sum properties. Location: \texttt{Theorem2/DiscreteCvM.lean}, lemma \texttt{abel\_sum\_identity}
\end{proof}

\subsection{Lag Regroup Identity (IMPLEMENTED in Lean)}

\begin{theorem}[Lag Regroup Identity]
\label{thm:lag-regroup}
For double sums over covariance structures with lag function:
\begin{equation}
\sum_{i=0}^{n-1} \sum_{j=0}^{n-1} a_i a_j \gamma(j-i) = \sum_{m=-(n-1)}^{n-1} \gamma(m) \cdot \text{coeff\_at\_lag}(a, m, n)
\end{equation}
where $\text{coeff\_at\_lag}(a, m, n) = \sum_{i: \text{valid}} a_i a_{i+m}$ has at most $n - |m|$ terms.
\end{theorem}

\begin{proof}[Formalization Status]
Location: \texttt{Proofs/LagRegroupProof.lean}. The proof uses \texttt{Finset.sum\_bij} for bijective reindexing, establishing the combinatorial identity. Full de-axiomatization status: \textbf{est.~4--8 weeks}
\end{proof}

\subsection{Davydov-Covariance Bound (AXIOMATIZED in Lean)}

\begin{theorem}[Davydov-Covariance Bound]
\label{thm:davydov-cov}
For the kernel-weighted process with lag-$d$ covariance $\gamma_d(y,z)$:
\begin{equation}
|\gamma_d(y,z)| \leq 4 \cdot \alpha(d)^{\frac{\delta}{2+\delta}}
\end{equation}
where $\alpha(d)$ are strong mixing coefficients and $\delta > 0$.
\end{theorem}

\begin{proof}[De-axiomatization Path]
Location: \texttt{Lemma1/Mixing.lean} and \texttt{Lemma1/Summability.lean}. The \textbf{critical path} involves:
\begin{enumerate}
    \item \texttt{Lemma1/Mixing.lean}: $\sigma$-algebra inclusion lemmas (sigma\_algebra\_inclusion, indicator\_measurable)
    \item \texttt{Lemma1/Mixing.lean}: Full Rio/Davydov inequality with exponent $\frac{\delta}{2+\delta}$ (est.\ 6--12 months for Mathlib $L^p$ infrastructure)
    \item \texttt{Lemma1/Summability.lean}: Comparison test for summability (est.\ 2--4 weeks after Davydov)
\end{enumerate}
\end{proof}

\begin{remark}[Rio Exponent]
The exponent $\frac{\delta}{2+\delta}$ comes from optimizing the Davydov bound in $L^{2+\delta}$ spaces (Rio 2013, Theorem~1.1). For centered indicators with $p = 2+\delta$, this gives sharper decay than the classical $L^\infty$ bound with exponent $1$.
\end{remark}

% ============================================
\section{Conclusion}
\label{sec:conclusion}
% ============================================

We have established the mathematical framework for Spatial CvM testing with fixed bandwidth kernels. Key innovations include:

\begin{itemize}
    \item Non-degenerate asymptotics via fixed $h > 0$
    \item Lag regroup identity for efficient computation (formalized in \texttt{Proofs/LagRegroupProof.lean})
    \item Davydov-Rio covariance bounds with exponent $\frac{\delta}{2+\delta}$ (axiomatized)
    \item Abel summation for discrete approximations (proved in \texttt{Theorem2/DiscreteCvM.lean})
\end{itemize}

The Lean 4 formalization tracks:
\begin{itemize}
    \item \textbf{PROVED:} Abel summation, basic kernel properties, lag regroup framework
    \item \textbf{AXIOMATIZED:} Davydov inequality (Rio 2013), Prokhorov's theorem, Mercer decomposition
    \item \textbf{ESTIMATED COMPLETION:} 8--16 months (with Rio/Davydov framework), down from 2--4 years
\end{itemize}

The structured formalization provides a rigorous roadmap with explicit mathematical foundations (Davydov 1968; Rio 2013) and clear decomposition of what requires Mathlib infrastructure versus what is already available.

% ============================================
% References
% ============================================

\begin{thebibliography}{99}

\bibitem[Davydov(1968)]{davydov1968}
Davydov, Y. A. (1968).
Convergence of distributions generated by stationary stochastic processes.
\textit{Theory of Probability \& Its Applications}, 13(4), 691--696.

\bibitem[Mercer(1909)]{mercer1909}
Mercer, J. (1909).
Functions of positive and negative type, and their connection with the theory of integral equations.
\textit{Philosophical Transactions of the Royal Society of London}, 209(441-458), 415--446.

\bibitem[Rio(2013)]{rio2013}
Rio, E. (2013).
\textit{Inequalities and limit theorems for weakly dependent sequences}.
Habilitation thesis, Universit\'{e} de Toulouse.\\
\url{https://www.ceremade.dauphine.fr/~ede/SH2013.pdf}

\bibitem[Skorokhod(1956)]{skorokhod1956}
Skorokhod, A. V. (1956).
Limit theorems for stochastic processes.
\textit{Theory of Probability \& Its Applications}, 1(3), 261--290.

\bibitem[El~Machkouri et al.(2013)]{elmachkouri2013}
El Machkouri, M., Voln\'{y}, D., \& Wu, W. B. (2013).
A central limit theorem for stationary random fields.
\textit{Stochastic Processes and their Applications}, 123(1), 1--14.

\end{thebibliography}

\end{document}
